\documentclass[11pt]{article}
\usepackage{fontspec,amsmath,amssymb,amsthm,hyperref}
\usepackage[margin=1in]{geometry}
\setlength{\emergencystretch}{3em}
\title{Optimizing friction and the shifted contraction exponent}
\author{LangevinLean mathematical exposition}
\date{}
\begin{document}\maketitle
\noindent Source module: \texttt{KineticLangevin/Friction.lean}. Compile this proof document with XeLaTeX or LuaLaTeX. The accompanying HTML pairs the mathematical argument with the exact Lean statements and records the scope of the formalization.


The general kinetic theorem permits every shift \(0\le a<\gamma/2\). Its exponential coefficient is \(k(a)=\max\{a^2-m,M-(\gamma-a)^2\}/(\gamma-2a)\); a negative coefficient is a contraction rate. This module proves the three scalar optimization regimes in Table 1 of the supplied manuscript. It establishes both the value at the proposed shift and a lower bound valid for every admissible competing shift. Singular square-root and denominator boundaries are explicit. These are algebraic corollaries of the kinetic estimate, separate from construction of the Brownian process and its invariant law.


\section{Center the shift interval}


Write \(d=\gamma-2a>0\). The two branches of the coefficient become

\[k(a)=-\frac\gamma2+\max\left\{\frac{d^2+\gamma^2-4m}{4d},\frac{4M-\gamma^2-d^2}{4d}\right\}.\]

The definition records the original coefficient exactly. For the centered identity, Lean proves each rational identity by expressing its numerator as \((-\gamma/2)d\) plus a quarter of the new numerator, distributing division, and cancelling the nonzero \(d\). Positivity allows the maximum to commute with division. Translating both branches by the same scalar then extracts \(-\gamma/2\). The denominator sign matters; a mere nonzero premise would not justify the maximum identity.



\noindent Lean declarations: \texttt{contractionExponent}, \texttt{contractionExponent\_centered}.

\section{Prove optimality of zero shift at low friction}


If \(\gamma^2\le m+M\) and \(m\le M\), the optimum is \(a=0\), with coefficient \((M-\gamma^2)/\gamma\). For any admissible \(a\ge0\), compare only the second branch to this value. After multiplying by the two positive denominators, their difference is \(a(2M-\gamma a)\), which is nonnegative under the regime assumptions. The full maximum is at least its second branch. Lean formalizes those two comparisons with the positive-denominator division lemmas and \texttt{nlinarith}. At zero, \(M-\gamma^2\ge-m\), so the maximum is exactly the second endpoint. In particular, \(\gamma^2\le M\) gives a nonnegative best exponent, explaining the table's absence of a contraction rate in that range.



\noindent Lean declarations: \texttt{contractionExponent\_low\_optimal}, \texttt{contractionExponent\_low\_zero}.

\section{Balance the endpoints in the middle regime}


Set \(D=\sqrt{2(m+M)-\gamma^2}\) and \(a=(\gamma-D)/2\). If \(\gamma^2\ge m+M\) and the radicand is strictly positive, then \(0<D\le\gamma\); hence the shift is admissible. Substituting \(D^2=2(m+M)-\gamma^2\) shows that the two curvature branches agree. The resulting coefficient is

\[k(a)=-\frac\gamma2+\frac{M-m}{2D}.\]

Lean proves the root sign and square first, establishes the denominator identity \(\gamma-2a=D\), rewrites the equal branches, and clears the positive denominator. It never evaluates the formula with \(D=0\).



\noindent Lean declarations: \texttt{contractionExponent\_middle\_value}.

\section{Check every competing shift at the middle corner}


Assume additionally \(\gamma^2\le3m+M\), and let \(d=\gamma-2a>0\) describe any competing shift. If \(d\le D\), compare the second branch with the balanced value. Clearing positive denominators factors its error as a positive multiple of \((D-d)(D^2+Dd+2(M-m))\). If \(d\ge D\), compare the first branch instead; the corresponding factor is \((d-D)(D(d-D)+2D^2-2(M-m))\). The last coefficient is nonnegative because \(\gamma^2\le3m+M\). Lean splits precisely on \(d\le D\), proves these products nonnegative, and uses the square identity for \(D\), also multiplied by \(D\), to close the polynomial comparisons. This proves global optimality, including the nondifferentiable corner of the maximum, without assuming a derivative there.



\noindent Lean declarations: \texttt{contractionExponent\_middle\_optimal}.

\section{Select the smaller characteristic root at high friction}


For \(\gamma^2\ge3m+M\), let \(D=\sqrt{\gamma^2-4m}>0\) and \(a=(\gamma-D)/2\). Nonnegative \(m\) gives \(D\le\gamma\), so the shift is nonnegative and its coordinate denominator is \(D>0\). The identity \(a(\gamma-a)=m\) makes the first branch equal \(-a\), and the high-friction condition places the second branch below it. Therefore the rate is \(a=(\gamma-\sqrt{\gamma^2-4m})/2\).



For any competing shift, the first branch alone bounds the maximum from below. The inequality that it is at least \(-a\) is exactly the nonnegative square \((\gamma-2a'-D)^2\ge0\) after division by the positive competing denominator. Lean supplies that square to \texttt{nlinarith}, then uses the maximum comparison. The strict discriminant excludes the equal-curvature critical point where the selected transformation becomes singular; a separate limiting argument would be required at that point.



\noindent Lean declarations: \texttt{contractionExponent\_high\_value}, \texttt{contractionExponent\_high\_optimal}.

\end{document}
